\section{Conclusion}\label{sec:conclusion}

We asked whether a theory of everything could ever be known to be final, and argued that
it could not.
\begin{enumerate}[label=(C\arabic*)]
  \item No method of any kind can certify a theory as final: if a method certifies a
        model, deterministic or stochastic, at a record, a computable world that produces
        that record makes the certification wrong (\cref{thm:k1-certify}).
  \item No usable method is guaranteed, even with any fixed probability, to settle on a
        final theory. Whatever has been observed so far, there is a possible world that
        fits all of it in which a deterministic usable method never settles on the right
        theory (\cref{thm:ud-gold}), and for a randomized one and each $p>0$, there is a world in which it settles with
        probability less than $p$ (\cref{ud:rem-randomized}). This part rests on the second premise that usable methods
        are computable (\cref{prem:effective}), and for deterministic methods only on its
        weaker form (\cref{thm:ud-gold-rel}).
  \item Even an ideal agent cannot separate formulations with identical predictions,
        and cannot tell apart with certainty chance hypotheses that agree on which events
        have probability zero (\cref{thm:t5-permanent}).
\end{enumerate}
As a corollary (C$'$), simplicity narrows the field but cannot decide among the survivors
that matter. The data fix which models fit; which of them is shortest depends on the
reference machine. Rivals that fit equally well, are nearly as short and disagree about
the unobserved always survive, and a change of machine that moves no description length by more than about a rival's gap makes it the
unique shortest (\cref{thm:con-gap,con:cor:promote}). The universal mixture still predicts
within $\K(\nu)$ bits of every model $\nu$ (\cref{thm:pred-relative}). Every choice among them is a
constraint we impose (\cref{lem:bias}).

The position is agnostic in Huxley's sense. It does not deny that the world has a true
description, or that physicists might one day write one down; it denies that anyone
could know, from the records, that they had. Nor does it counsel giving up on unification. Compact theories remain
worth seeking, because they compress what is known and commit to new, testable
predictions (\cref{sec:predict}). What the argument removes is the goal of a final theory known to be final, and with it the hope that simplicity, beauty
or unification could certify one.

Put plainly: we can keep finding better theories of everything, but we can never know,
from the records, that we have found the last one.
